\section{From Threshold Straddling to Literal Affine Crossings}
\label{sec:affine}

The Boolean square records which threshold decisions differ, but it does not yet say what numerical event lies behind a changed bit. The next formal layer reconnects categorical displacement to rational support masses.

\subsection{Endpoint straddling}

For a rational threshold \(c\), define
\[
  \operatorname{Straddle}(c;x,y)
\]
to mean that exactly one endpoint lies on or above the threshold:
\[
  (c\le x\land\neg(c\le y))
  \ \lor\ 
  (c\le y\land\neg(c\le x)).
\]
For rational numbers this is the expected endpoint pattern
\[
  x<c\le y
  \quad\text{or}\quad
  y<c\le x.
\]

\begin{theorem}[Threshold decision versus endpoint straddling]
For rational \(c,x,y\),
\[
  [c\le x]\neq[c\le y]
  \iff
  \operatorname{Straddle}(c;x,y).
\]
\statusverified
\end{theorem}

At belief level, the total wall count decomposes into the one-dimensional wall counts of the positive and negative support coordinates. Lean therefore verifies the exact three-way correspondence:
\[
\begin{array}{ccl}
\dthr=0 &\iff& \text{neither support coordinate straddles},\\
\dthr=1 &\iff& \text{exactly one support coordinate straddles},\\
\dthr=2 &\iff& \text{both support coordinates straddle}.
\end{array}
\]
In particular, every genuine categorical belief change has a numerical endpoint witness on at least one signed support axis.

\subsection{Affine interpolation}

Endpoint straddling still describes only two snapshots. To expose an actual crossing event, the formalization introduces the directed rational affine path
\[
  \gamma_{x,y}(t)=x+t(y-x),
  \qquad t\in\mathbb{Q}.
\]
It satisfies
\[
  \gamma_{x,y}(0)=x,
  \qquad
  \gamma_{x,y}(1)=y.
\]

For the increasing orientation \(x<c\le y\), the explicit threshold parameter is
\[
  t_+=\frac{c-x}{y-x}.
\]
For the decreasing orientation \(y<c\le x\), the development uses the equivalent directed parameter
\[
  t_- =\frac{x-c}{x-y}.
\]
Lean Core rational arithmetic is sufficient to prove that the relevant parameter lies in the unit interval and hits the threshold exactly.

\begin{theorem}[Affine unit-interval threshold crossing]
If rational endpoints \(x,y\) straddle rational threshold \(c\), then
\[
  \exists t\in\mathbb{Q}:
  0\le t\le 1
  \quad\land\quad
  \gamma_{x,y}(t)=c.
\]
\statusverified
\end{theorem}

\begin{proof}[Proof sketch]
Straddling implies distinct endpoints. Solve \(x+t(y-x)=c\) for \(t=(c-x)/(y-x)\); the straddling inequalities give \(0\le t\le1\), accounting for the sign of the denominator. Substitution verifies the hit. A hit can be at an endpoint; the theorem does not require \(0<t<1\).
\end{proof}
No abstract topology is used, and there is no Mathlib dependency.

\begin{corollary}[Belief change has a literal affine wall hit]
If admissible conditionalization changes the complete FDE value of \(\Bel_i\varphi\), then at least one of the two chosen interpolations of its endpoint support masses
\[
  t\mapsto \Ppos_t(\varphi),
  \qquad
  t\mapsto \Pneg_t(\varphi)
\]
hits \(c_i\) at some rational \(t\in[0,1]\). Here \(\Ppos_t=(1-t)\Ppos_0+t\Ppos_1\), and likewise for negative support; this notation does not yet denote evaluation in an intermediate model.
\statusverified
\end{corollary}

For a two-wall transition, both signed support coordinates have crossing witnesses, though they need not cross at the same parameter. This asks an order question about the chosen interpolation: which wall is hit first?

\paragraph{Support-plane picture.}
The two masses define a point
\[
  (\Ppos,\Pneg)\in\mathbb{Q}^2.
\]
The threshold lines \(\Ppos=c\) and \(\Pneg=c\) partition the plane into the four FDE regions.

\begin{center}
\begin{tikzpicture}[scale=1.25]
  \draw[->] (0,0) -- (5,0) node[right] {$\Ppos$};
  \draw[->] (0,0) -- (0,5) node[above] {$\Pneg$};
  \draw[dashed] (2.4,0) -- (2.4,5) node[above] {$\Ppos=c$};
  \draw[dashed] (0,2.4) -- (5,2.4) node[right] {$\Pneg=c$};
  \node at (1.2,1.2) {$\Nval$};
  \node at (3.7,1.2) {$\Tval$};
  \node at (1.2,3.7) {$\Fval$};
  \node at (3.7,3.7) {$\Bval$};
  \fill (2.4,2.4) circle (1.5pt);
  \node[above right] at (2.4,2.4) {$(c,c)$};
\end{tikzpicture}
\end{center}

The threshold square is the four-region categorical observation of the signed support plane. These four regions must not be confused with the four evidence masses \((t,b,n,f)\) in a probability simplex. For integrity-certified models the relevant support points lie in the unit square; the algebraic lemmas themselves permit arbitrary rational coordinates.
